\documentclass[UTF8,fontset=none,11pt]{ctexart}
\usepackage[a4paper,margin=2.05cm]{geometry}
\usepackage{amsmath,amssymb}
\usepackage{enumitem}
\usepackage[table]{xcolor}
\usepackage{booktabs}
\usepackage{array}
\usepackage{longtable}
\usepackage{tikz}
\usepackage{graphicx}
\usepackage{hyperref}
\usetikzlibrary{positioning,arrows.meta,shapes.geometric,fit,backgrounds,calc}

\setCJKmainfont[AutoFakeSlant=0.18]{PingFang SC}
\setCJKsansfont[AutoFakeSlant=0.18]{PingFang SC}
\setCJKmonofont[AutoFakeSlant=0.18]{PingFang SC}

\hypersetup{
  colorlinks=true,
  linkcolor=blue!55!black,
  urlcolor=blue!55!black,
  citecolor=blue!55!black,
  pdftitle={Target Code Generation 中文学习笔记},
  pdfauthor={Trae AI}
}

\setlength{\parindent}{0pt}
\setlength{\parskip}{0.55em}
\setlist[itemize]{leftmargin=1.8em,itemsep=0.22em,topsep=0.25em}
\setlist[enumerate]{leftmargin=2.0em,itemsep=0.28em,topsep=0.3em}
\renewcommand{\arraystretch}{1.18}
\setlength{\emergencystretch}{3em}
\sloppy

\definecolor{defrule}{RGB}{25,92,140}
\definecolor{noterule}{RGB}{120,80,15}
\definecolor{warnrule}{RGB}{160,50,45}
\definecolor{softblue}{RGB}{235,246,255}
\definecolor{softyellow}{RGB}{255,249,230}
\definecolor{graytext}{RGB}{95,95,95}
\definecolor{GRAYTEXT}{RGB}{95,95,95}
\newcolumntype{L}[1]{>{\raggedright\arraybackslash}p{#1}}

\newenvironment{defbox}[1][]{%
  \par\medskip\noindent{\color{defrule}\rule{\linewidth}{1.1pt}}\par\nopagebreak
  \noindent{\bfseries\color{defrule}#1}\par\nopagebreak\smallskip}%
  {\par\nopagebreak\smallskip\noindent{\color{defrule}\rule{\linewidth}{0.4pt}}\par\medskip}
\newenvironment{notebox}[1][]{%
  \par\medskip\noindent{\color{noterule}\rule{\linewidth}{1.1pt}}\par\nopagebreak
  \noindent{\bfseries\color{noterule}#1}\par\nopagebreak\smallskip}%
  {\par\nopagebreak\smallskip\noindent{\color{noterule}\rule{\linewidth}{0.4pt}}\par\medskip}
\newenvironment{warnbox}[1][]{%
  \par\medskip\noindent{\color{warnrule}\rule{\linewidth}{1.1pt}}\par\nopagebreak
  \noindent{\bfseries\color{warnrule}#1}\par\nopagebreak\smallskip}%
  {\par\nopagebreak\smallskip\noindent{\color{warnrule}\rule{\linewidth}{0.4pt}}\par\medskip}

\newcommand{\pg}[1]{\texorpdfstring{{\small\color{graytext}\,[p#1]}}{ [p#1]}}
\newcommand{\IR}{\ensuremath{\mathrm{IR}}}
\newcommand{\TAC}{\ensuremath{\mathrm{TAC}}}
\newcommand{\CFG}{\ensuremath{\mathrm{CFG}}}
\newcommand{\RIG}{\ensuremath{\mathrm{RIG}}}
\newcommand{\AR}{\ensuremath{\mathrm{AR}}}

\title{\bfseries Lecture 9: Target Code Generation\\中文学习笔记}
\author{基于课件 \texttt{9.Target\_Code\_Generation.pdf} 整理}
\date{\today}

\begin{document}
\maketitle
\tableofcontents
\newpage

\section*{阅读导引}
\addcontentsline{toc}{section}{阅读导引}

本讲进入编译器后端的目标代码生成阶段。前面已经得到优化后的中间表示 IR；现在要把 IR 变成目标机器能执行的汇编或机器码。目标代码生成的核心矛盾是：IR 中临时变量很多、操作很抽象，而真实机器寄存器有限、指令集有具体限制、函数调用还必须遵守运行时栈和调用约定。

\begin{notebox}[本讲主线]
目标代码生成要解决三件事：选什么指令、哪些值放寄存器、指令按什么顺序执行。为了说明这些问题，课件先用栈式机器和 MIPS 建模，再讲运行时栈和调用约定，最后讨论对象分派和机器相关优化。
\end{notebox}

\section{编译阶段与目标代码生成任务 \pg{1--6}}

\subsection{目标代码生成的位置 \pg{1--3}}

目标代码生成（target code generation）位于优化后的中间代码和目标机器代码之间：
\[
\text{IR}\to \text{Code Optimizer}\to \text{Code Generator}\to \text{target-machine code}.
\]

\begin{defbox}[Target Code Generation]
目标代码生成是把语法分析、语义分析和优化后得到的中间代码转换为目标代码的阶段。目标代码可以是汇编代码，也可以是机器码。
\end{defbox}

这一阶段要考虑：

\begin{itemize}
  \item 如何让目标代码更短。
  \item 如何充分使用寄存器，减少内存访问。
  \item 如何利用目标机器指令系统的特点。
\end{itemize}

输入通常包括源程序 IR 和符号表；输出是目标程序。

\subsection{三个主要任务 \pg{4--6}}

\begin{longtable}{L{0.23\linewidth}L{0.33\linewidth}L{0.34\linewidth}}
\toprule
任务 & 含义 & 难点 \\
\midrule
Instruction selection & 选择合适的目标机器指令实现 IR 语句 & 同一 IR 操作可能有多种目标指令实现 \\
Register allocation and assignment & 决定哪些值保存在什么寄存器中 & IR 临时变量近似无限，物理寄存器很少 \\
Instruction ordering / scheduling & 决定指令执行顺序 & 顺序影响寄存器需求、流水线等待和并行性 \\
\bottomrule
\end{longtable}

指令选择例子：

\begin{verbatim}
a = b + c;
d = a + e;
\end{verbatim}

可翻译为：

\begin{verbatim}
LD  R0, b
ADD R0, R0, c
ST  a, R0
LD  R0, a
ADD R0, R0, e
ST  d, R0
\end{verbatim}

如果后端能发现 \(a\) 的值仍在 \(R0\)，就可能避免重新 `LD R0, a`。这正是寄存器分配和指令调度会影响代码质量的原因。

寄存器分配和最优指令顺序选择在一般情况下很难，课件指出它们与 NP 问题相关；实践中通常使用启发式算法，而不是追求严格最优。

\section{栈式计算机、累加器与 MIPS 基础 \pg{7--13}}

\subsection{栈式计算机 \pg{7}}

\begin{defbox}[Stack Machine]
栈式计算机是一种简单求值模型：没有显式变量或寄存器，中间结果都放在一个值栈中。每条指令从栈顶取操作数、弹出操作数、完成运算，并把结果压回栈顶。
\end{defbox}

通俗理解：表达式求值像用一摞盘子做计算，最新产生的值总放在最上面，下一条指令从最上面拿操作数。

它的优点是模型简单；缺点是栈顶访问非常频繁，导致内存读写多。

\subsection{用累加器优化栈式机器 \pg{8--9}}

普通 `add` 指令需要三次内存操作：读两个操作数，再写回一个结果。课件提出优化：把栈顶值保存在一个寄存器中，这个寄存器称为 accumulator（累加器）。

此时加法变成：
\[
acc \leftarrow acc + top\_of\_stack.
\]
这样只需要访问一次内存，即读取栈中另一个操作数。

对表达式 \(3+7+5\)，栈式求值会依次把操作数压栈/放入累加器，再做两次加法，最终结果保存在累加器中。

\subsection{从栈式机器到 MIPS \pg{10--13}}

课件用 MIPS 模拟带累加器的栈式机器：

\begin{itemize}
  \item ACC 存在 MIPS 寄存器 \texttt{\$t0}。
  \item 栈存在内存中。
  \item \texttt{\$sp} 保存栈上下一可用位置地址。
  \item 栈顶位于 \texttt{\$sp + 4}。
  \item 栈向低地址增长，这是 MIPS 的标准约定。
\end{itemize}

\begin{defbox}[MIPS]
MIPS 是典型 RISC 架构。算术指令只对寄存器操作；若要使用内存中的值，必须先用 load 指令读入寄存器，计算后再用 store 写回内存。
\end{defbox}

MIPS 特点：

\begin{itemize}
  \item 32 个 32-bit 通用寄存器。
  \item 本讲主要用 \texttt{\$t0} 作为 ACC，\texttt{\$sp} 作为栈指针，\texttt{\$t1} 作为临时寄存器。
  \item \texttt{\$zero} 恒为 0。
\end{itemize}

常见指令：

\begin{longtable}{L{0.27\linewidth}L{0.63\linewidth}}
\toprule
指令 & 含义 \\
\midrule
\texttt{lw reg1 offset(reg2)} & 从地址 \texttt{reg2+offset} 读取 32-bit word 到 \texttt{reg1} \\
\texttt{add reg1 reg2 reg3} & \texttt{reg1 <- reg2 + reg3} \\
\texttt{sw reg1 offset(reg2)} & 把 \texttt{reg1} 的 32-bit word 写到地址 \texttt{reg2+offset} \\
\texttt{addiu reg1 reg2 imm} & \texttt{reg1 <- reg2 + imm}，不检查溢出 \\
\texttt{li reg imm} & 把立即数装入寄存器 \\
\bottomrule
\end{longtable}

\section{寄存器保存恢复、栈操作和表达式代码生成 \pg{14--20}}

\subsection{为什么要保存/恢复寄存器 \pg{14--15}}

IR 中可以有无限临时变量，但真实机器寄存器有限，因此必须复用寄存器。复用带来问题：如果一个寄存器中旧值之后还要用，而我们要把新值写进去，就必须先保存旧值，之后再恢复。

例如生成 \(E\to E_1+E_2\)：

\begin{itemize}
  \item 先生成 \(E_1\)，结果在 \texttt{\$t0}。
  \item 再生成 \(E_2\)，也会把结果写到 \texttt{\$t0}，覆盖 \(E_1\) 的结果。
  \item 因此在生成 \(E_2\) 前，要把 \(E_1\) 的 \texttt{\$t0} 保存到栈中。
\end{itemize}

保存 \texttt{\$t0}：
\begin{verbatim}
sw    $t0, 0($sp)
addiu $sp, $sp, -4
\end{verbatim}

恢复到 \texttt{\$t0}：
\begin{verbatim}
lw    $t0, 4($sp)
addiu $sp, $sp, 4
\end{verbatim}

这里栈向低地址增长；压栈后 \texttt{\$sp} 下降，弹栈时 \texttt{\$sp} 上升。

\subsection{栈操作 \pg{16--18}}

压入多个寄存器时，通常先移动 \texttt{\$sp} 为新数据留空间，再保存：

\begin{verbatim}
addiu $sp, $sp, -8
sw    $t1, 4($sp)
sw    $t2, 0($sp)
\end{verbatim}

弹出时只需要把 \texttt{\$sp} 向高地址移动。被弹出的数据仍在内存里，但位于栈指针之外，按约定视为无效。

\subsection{MIPS 翻译例子：\texorpdfstring{\(7+5\)}{7+5} \pg{17}}

栈式机器代码：

\begin{verbatim}
acc <- 7
push acc
acc <- 5
acc <- acc + stack_top
pop
\end{verbatim}

MIPS 翻译：

\begin{verbatim}
li    $t0, 7
sw    $t0, 0($sp)
addiu $sp, $sp, -4
li    $t0, 5
lw    $t1, 4($sp)
add   $t0, $t0, $t1
addiu $sp, $sp, 4
\end{verbatim}

执行后，\texttt{\$t0} 中为 12，\texttt{\$t1} 中为 7，\texttt{\$sp} 回到压栈前位置。内存中原来保存 7 的位置可能仍有 7，但已经不属于有效栈内容。

\subsection{表达式代码生成策略 \pg{19}}

定义代码生成函数 \texttt{cgen(e)}：

\begin{itemize}
  \item 生成表达式 \(e\) 的 MIPS 代码。
  \item 执行后 \(e\) 的值保存在 \texttt{\$t0}。
  \item 保持 \texttt{\$sp} 和栈内容一致。
\end{itemize}

对加法：

\begin{verbatim}
cgen(e1 + e2):
  cgen(e1)              # e1 result in $t0
  sw    $t0, 0($sp)
  addiu $sp, $sp, -4
  cgen(e2)              # e2 overwrites $t0
  lw    $t1, 4($sp)
  addiu $sp, $sp, 4
  add   $t0, $t1, $t0
\end{verbatim}

如果只用 \texttt{\$t1} 暂存 \(e_1\)，对简单表达式可少用栈：

\begin{verbatim}
cgen(e1)
move $t1, $t0
cgen(e2)
add  $t0, $t1, $t0
\end{verbatim}

但若 \(e_2\) 自身也很复杂，例如 \(3+(7+5)\)，\texttt{\$t1} 可能被内部计算需要，仍需更系统的寄存器分配。

\subsection{条件语句代码生成 \pg{20}}

新增控制流指令：

\begin{itemize}
  \item \texttt{beq reg1 reg2 label}：若两个寄存器相等，跳到 label。
  \item \texttt{b label}：无条件跳转。
\end{itemize}

对：
\[
if\ e_1==e_2\ then\ e_3\ else\ e_4
\]
代码生成：

\begin{verbatim}
cgen(e1)
sw    $t0, 0($sp)
addiu $sp, $sp, -4
cgen(e2)
lw    $t1, 4($sp)
addiu $sp, $sp, 4
beq   $t0, $t1, true_branch
false_branch:
  cgen(e4)
  b end_if
true_branch:
  cgen(e3)
end_if:
\end{verbatim}

基本思路是：比较前先保存 \(e_1\)，再计算 \(e_2\)，比较后选择 true/false 分支。

\section{运行时内存布局、活动记录和调用约定 \pg{21--30}}

\subsection{运行时内存区域 \pg{21}}

一个程序运行时常见内存布局：

\begin{longtable}{L{0.18\linewidth}L{0.72\linewidth}}
\toprule
区域 & 作用 \\
\midrule
Code & 保存目标代码，生成后大小通常编译期确定 \\
Global/static & 保存全局常量、静态对象等编译期大小已知的数据 \\
Heap & 管理长生命周期动态数据，如对象和动态分配数组 \\
Stack & 保存动态过程调用信息，尤其是活动记录、局部变量、临时值 \\
\bottomrule
\end{longtable}

低地址通常放 code/data，heap 向高地址增长，stack 从高地址向低地址增长，中间是 free memory。

\subsection{Activation 和 Activation Record \pg{22--24}}

\begin{defbox}[Activation]
一次过程或函数执行称为一个活动。活动从过程体开始执行时开始，过程结束后把控制权返回到调用点之后。
\end{defbox}

\begin{defbox}[Activation Record]
活动记录 AR 是管理某次过程执行所需信息的内存块，也称 stack frame 或 frame。调用栈中保存的就是一系列活动记录。
\end{defbox}

AR 栈管理：

\begin{itemize}
  \item 函数入口：在栈顶分配新的 AR。
  \item 函数返回：从栈顶移除该 AR。
  \item \texttt{\$sp} 指向栈顶，移动 \texttt{\$sp} 即可分配/释放空间。
  \item \texttt{\$fp} 保存当前 frame 的基地址，局部变量和参数可用相对 \texttt{\$fp} 的偏移访问。
\end{itemize}

AR 内容通常包括：

\begin{itemize}
  \item 临时变量。
  \item 局部变量。
  \item 保存的调用者/被调用者寄存器值。
  \item 调用者返回地址，也就是调用后下一条指令地址。
  \item 调用者的旧 \texttt{\$fp}。
  \item 参数。
  \item 返回值。
\end{itemize}

\subsection{Caller/Callee Conventions \pg{25--27}}

函数调用会覆盖寄存器。问题是：由调用者保存，还是被调用者保存？

\begin{itemize}
  \item 调用者知道自己哪些寄存器值之后还要用。
  \item 被调用者知道自己会用哪些寄存器。
  \item 但二者不知道对方实现细节，因此需要约定。
\end{itemize}

课件给出合作方案：

\begin{longtable}{L{0.22\linewidth}L{0.27\linewidth}L{0.41\linewidth}}
\toprule
类别 & 寄存器 & 责任 \\
\midrule
Caller-saved & \texttt{\$t0-\$t9}, \texttt{\$a0-\$a3}, \texttt{\$v0-\$v1} & 被调用者可以修改；调用者若关心这些值，调用前自己保存，调用后恢复 \\
Callee-saved & \texttt{\$s0-\$s7}, \texttt{\$ra} & 调用者假设这些值调用后不变；被调用者若使用，必须保存和恢复 \\
\bottomrule
\end{longtable}

\subsection{详细调用步骤 \pg{28--30}}

调用者：

\begin{enumerate}
  \item 为需要保存的 caller-saved 寄存器在栈上留空间并保存。
  \item 参数从右到左逐个压栈。
  \item 跳转到函数入口，并保存下一条指令地址到 \texttt{\$ra}。
\end{enumerate}

被调用者：

\begin{enumerate}
  \item 在栈上留空间，保存旧 \texttt{\$fp} 和 \texttt{\$ra}。
  \item 设置自己的 \texttt{\$fp} 为当前 frame 的基地址。
  \item 为局部变量和临时变量分配栈帧空间。
  \item 执行函数体。
\end{enumerate}

返回时，被调用者：

\begin{enumerate}
  \item 将返回值放入 \texttt{\$v0}。
  \item 弹出局部变量、临时变量、保存的寄存器。
  \item 恢复 \texttt{\$fp} 和 \texttt{\$ra}。
  \item 执行 \texttt{jr \$ra} 跳回调用者。
\end{enumerate}

调用者恢复：

\begin{enumerate}
  \item 弹出参数。
  \item 恢复自己保存的 caller-saved 寄存器。
\end{enumerate}

新增指令：

\begin{itemize}
  \item \texttt{jal label}：跳到函数入口，并把下一条指令地址保存到 \texttt{\$ra}。
  \item \texttt{jr reg}：跳到寄存器中的地址，函数返回常用 \texttt{jr \$ra}。
\end{itemize}

\section{变量访问与帧指针偏移 \pg{31--33}}

函数的参数、局部变量和临时变量都在 AR 中。问题是：如果用 \texttt{\$sp} 相对地址访问变量，表达式求值期间不断压栈/弹栈会移动 \texttt{\$sp}，导致变量偏移不稳定。

解决方案是使用 \texttt{\$fp}：

\begin{itemize}
  \item \texttt{\$fp} 在当前函数执行期间保持不动。
  \item 参数、局部变量和临时变量都可用固定 offset 访问。
  \item 参数通常在 \texttt{\$fp} 的正偏移处。
  \item 局部变量和临时变量通常在 \texttt{\$fp} 的负偏移处。
\end{itemize}

课件示例：

\begin{verbatim}
def f(x, y) = e
x: +8($fp)
y: +12($fp)
first local variable: -4($fp)
\end{verbatim}

对嵌套调用，例如 \texttt{fun1} 调用 \texttt{fun2}，栈上会同时存在两个 AR；每个函数都有自己的 \texttt{\$fp}，因此相同 offset 在不同 AR 中指向不同函数的变量。

\section{面向对象代码生成：对象布局与分派 \pg{34--39}}

\subsection{对象布局 \pg{34}}

对象类似 C 结构体：

\begin{itemize}
  \item 对象在内存中连续存放。
  \item 每个成员变量位于对象内固定 offset。
  \item 与结构体不同，对象还有成员方法。
\end{itemize}

成员访问可翻译为 offset 访问。例如：

\begin{verbatim}
obj.x = 10
\end{verbatim}

可翻译为“把 10 写到 \texttt{obj + x\_offset}”。

\subsection{非虚方法、虚方法与分派 \pg{35--36}}

\begin{defbox}[Static Dispatch]
静态分派是在编译期确定调用目标。非虚方法不能被覆盖，调用目标由静态引用类型决定，因此编译器可把函数地址硬编码到目标代码中。
\end{defbox}

\begin{defbox}[Dynamic Dispatch]
动态分派是在运行时根据对象真实类型选择调用目标。虚方法可被子类覆盖，因此同一个调用表达式在运行时可能跳到不同函数。
\end{defbox}

例子：

\begin{verbatim}
Parent obj = new Child();
obj.nonvirtual();  // 调用 Parent::nonvirtual()
obj.virtual();     // 调用 Child::virtual()
\end{verbatim}

动态分派的实现方式：

\begin{itemize}
  \item 编译期为每个类生成 dispatch table，表中保存该类所有虚方法的目标地址。
  \item 运行时每个对象带一个 dispatch pointer，指向其真实类的 dispatch table。
  \item 调用虚方法时，通过 dispatch pointer 加方法 offset 取出函数地址，再跳转调用。
\end{itemize}

\subsection{典型对象布局与继承不变式 \pg{37--39}}

典型对象布局：

\begin{itemize}
  \item Class tag：用于动态类型检查。
  \item Dispatch ptr：指向 dispatch table。
  \item 成员变量，如 \(x,y,z\)。
\end{itemize}

继承中必须保持一个不变式：

\begin{defbox}[Offset Invariant]
某个成员变量或成员方法在父类及其所有子类中的 offset 必须一致。这样通过父类引用访问对象时，即使对象真实类型是子类，编译器仍可使用父类布局的 offset。
\end{defbox}

例如：

\begin{itemize}
  \item A1 有 \(x\)、虚方法 \(f1,f2\)。
  \item A2 继承 A1，新增 \(y\)，覆盖 \(f2\)。
  \item A3 继承 A2，新增 \(z\)，新增 \(f3\)。
\end{itemize}

布局方式：

\begin{itemize}
  \item 子类保留父类字段 offset，新增字段追加在后面。
  \item 子类覆盖虚方法时，dispatch table 中对应 offset 的函数地址替换成子类版本。
  \item 新增虚方法追加到 dispatch table 后面。
\end{itemize}

成员变量访问用引用类型对应的 offset；虚方法调用也用引用类型对应的 dispatch table offset。由于 offset invariant，真实对象是子类也能正确工作。

\section{机器相关优化 \pg{40--49}}

\subsection{机器相关优化总览 \pg{40}}

IR 优化后，编译器需要把优化后的 IR 转成目标语言。此时必须考虑具体架构特征：

\begin{itemize}
  \item 专用指令。
  \item 硬件流水线能力。
  \item 寄存器数量和使用约定。
\end{itemize}

典型机器相关优化：

\begin{itemize}
  \item 指令选择与调度。
  \item 寄存器分配。
  \item 窥孔优化。
\end{itemize}

\subsection{指令选择与指令成本 \pg{41--43}}

指令选择是在目标架构上为 IR 操作寻找高效汇编实现。它在 CISC 架构上尤其重要，因为一个 IR 操作可能有多种指令组合。

例如：

\begin{verbatim}
a = a + 1
\end{verbatim}

可以实现为：

\begin{verbatim}
MOV a, R0
ADD #1, R0
MOV R0, a
\end{verbatim}

也可以实现为：

\begin{verbatim}
INC a
\end{verbatim}

后者更短、更便宜。

课件给出指令成本公式：
\[
\text{Instruction cost}=1+\text{cost(source-mode)}+\text{cost(destination-mode)}.
\]

示例：

\begin{itemize}
  \item \texttt{a := b + c} 可用 \texttt{MOV b,R0; ADD c,R0; MOV R0,a}，成本约 6。
  \item 若允许直接内存目的操作，可用 \texttt{MOV b,a; ADD c,a}，成本也可能约 6。
  \item 对间接地址 \texttt{*R1,*R0} 的操作成本可能更高。
  \item \texttt{a := a + 1} 可从三条指令成本 6，降到 \texttt{ADD \#1,a} 成本 3，再降到 \texttt{INC a} 成本 2。
\end{itemize}

\subsection{指令调度 \pg{44}}

\begin{defbox}[Instruction Scheduling]
指令调度是重排目标指令顺序以减少执行时间的优化。它利用硬件流水线和多发射能力，尽量减少等待周期、避免寄存器溢出并改善局部性。
\end{defbox}

例子：

\begin{verbatim}
A = x * y;
B = A + 1;
C = y;
\end{verbatim}

若乘法延迟较长，执行 \(B=A+1\) 必须等待 \(A\) 完成；可改为：

\begin{verbatim}
A = x * y;
C = y;
B = A + 1;
\end{verbatim}

这样 \texttt{C=y} 可在等待乘法结果时执行，减少空转。

\subsection{寄存器分配 \pg{45--46}}

\begin{defbox}[Register Allocation]
寄存器分配是把变量和表达式结果映射到有限物理寄存器，并管理寄存器与内存之间数据转移的过程。
\end{defbox}

目标：

\begin{itemize}
  \item 频繁访问的变量尽量放入寄存器。
  \item 变量只在其活跃期间占用寄存器。
  \item 减少内存 load/store，因为内存访问代价高。
\end{itemize}

局部寄存器分配：

\begin{itemize}
  \item 逐基本块分配。
  \item 决策简单，但块之间映射可能不一致。
\end{itemize}

全局寄存器分配：

\begin{itemize}
  \item 跨基本块做整体决策。
  \item 变量到寄存器的映射可跨块保持一致。
  \item 需要考虑 CFG 结构。
\end{itemize}

常见算法：

\begin{itemize}
  \item 图着色分配器。
  \item 线性扫描分配器。
  \item 整数线性规划分配器。
\end{itemize}

\subsection{图着色与寄存器溢出 \pg{47--48}}

\begin{defbox}[Register Interference Graph]
寄存器干涉图 RIG 中，每个节点代表一个变量；若两个变量 live range 重叠，不能放入同一个寄存器，就在两个节点之间连边。
\end{defbox}

规则：

\begin{itemize}
  \item 没有边的两个变量可共享同一寄存器。
  \item 有边的两个变量不能共享同一寄存器。
  \item k 个颜色对应 k 个寄存器。
  \item 若 RIG 可 k 着色，则程序可用 k 个寄存器分配。
\end{itemize}

\begin{defbox}[Register Spilling]
若寄存器不够，编译器把某些变量溢出到内存中。被溢出的变量不长期占用寄存器，需要使用时从内存 load，用完后必要时 store 回内存。
\end{defbox}

k 着色判定是 NP-complete，因此 k 寄存器分配也是 NP-complete。实践中使用启发式多项式算法，例如 Chaitin 的图着色算法，GCC 后端中曾广泛使用类似思想。

溢出很慢，因为它引入额外内存读写。所以优秀的寄存器分配器对代码质量非常关键。

\subsection{窥孔优化 \pg{49}}

\begin{defbox}[Peephole Optimization]
窥孔优化是用一个小滑动窗口检查目标指令序列，并把窗口内指令替换为更短或更快序列的局部优化技术。
\end{defbox}

它通常有两种思路：

\begin{itemize}
  \item 通过良好的指令选择和寄存器分配直接生成高质量代码。
  \item 先生成朴素目标代码，再用窥孔优化局部改进。
\end{itemize}

常见变换：

\begin{itemize}
  \item 冗余指令消除。
  \item 控制流优化。
  \item 代数化简。
\end{itemize}

例子：jump to jump。

\begin{verbatim}
if a < b goto L1
...
L1: goto L2
\end{verbatim}

可改成：

\begin{verbatim}
if a < b goto L2
...
L1: goto L2
\end{verbatim}

条件跳转直接跳到最终目标，少执行一次无条件跳转。

\section{总结 \pg{50}}

代码可在不同层级优化：

\begin{itemize}
  \item 窥孔优化、局部优化、循环优化、全局优化。
  \item IR 层优化：局部/全局、公共子表达式消除、常量折叠和常量传播。
  \item 目标层优化：指令选择、指令调度、寄存器分配等。
\end{itemize}

优化之间相互影响：

\begin{itemize}
  \item 一个优化可能暴露另一个优化机会。
  \item 一个优化也可能隐藏或删除另一个优化机会。
  \item 最优优化组合仍是研究问题。
  \item 编译器通常把常一起使用的优化打包成等级，例如 \texttt{-O1}、\texttt{-O2}。
\end{itemize}

\section{练习题与解答}

\subsection*{题 1：目标代码生成的三个主要任务是什么？}

\textbf{解答：} 指令选择、寄存器分配与赋值、指令排序/调度。指令选择决定用哪些目标机器指令实现 IR；寄存器分配决定哪些值放入哪些寄存器；指令调度决定指令执行顺序，以减少等待并提高硬件利用率。

\subsection*{题 2：为什么栈式机器需要累加器优化？}

\textbf{解答：} 栈顶被频繁访问。普通栈式加法需要读两个操作数并写回结果，共三次内存访问。若把栈顶保存在累加器寄存器中，加法可变为 \(acc\leftarrow acc+top\_of\_stack\)，通常只需一次内存访问，减少开销。

\subsection*{题 3：执行课件中 \texorpdfstring{\(7+5\)}{7+5} 的 MIPS 代码后，\texttt{\$t0}、\texttt{\$t1} 和栈状态如何？}

\textbf{解答：} 执行后 \texttt{\$t0=12}，因为它保存 \(5+7\) 的结果；\texttt{\$t1=7}，因为它从栈顶恢复了先前保存的 7；\texttt{\$sp} 回到压入 7 之前的位置。原先保存 7 的内存单元可能仍有 7，但弹栈后不再属于有效栈内容。

\subsection*{题 4：写出 \texorpdfstring{\(e_1+e_2\)}{e1+e2} 的基本代码生成策略。}

\textbf{解答：}

\begin{verbatim}
cgen(e1)
sw    $t0, 0($sp)
addiu $sp, $sp, -4
cgen(e2)
lw    $t1, 4($sp)
addiu $sp, $sp, 4
add   $t0, $t1, $t0
\end{verbatim}

先计算 \(e_1\)，保存到栈；再计算 \(e_2\)，恢复 \(e_1\) 到 \texttt{\$t1}，最后相加。

\subsection*{题 5：为什么访问函数变量要用 \texttt{\$fp} 而不是 \texttt{\$sp}？}

\textbf{解答：} \texttt{\$sp} 会随着表达式求值、保存临时值、压栈/弹栈而变化，变量相对它的偏移不稳定。 \texttt{\$fp} 在当前函数执行期间保持固定，因此参数和局部变量可以用固定偏移访问。

\subsection*{题 6：Caller-saved 和 Callee-saved 的区别是什么？}

\textbf{解答：} Caller-saved 寄存器可被被调用者修改，调用者若需要保留它们的值，必须调用前保存、调用后恢复。Callee-saved 寄存器在调用后应保持不变，被调用者若使用它们，必须自己保存和恢复。

\subsection*{题 7：虚方法为什么需要动态分派？}

\textbf{解答：} 虚方法可被子类覆盖，调用目标取决于运行时对象真实类型，而不只是静态引用类型。因此编译期不能硬编码唯一目标，必须运行时通过对象的 dispatch pointer 查 dispatch table，找到真实类对应的函数地址。

\subsection*{题 8：继承中为什么要求成员变量和虚方法 offset 不变？}

\textbf{解答：} 父类引用可能指向子类对象。编译器根据引用类型生成固定 offset 的字段访问或虚方法表访问；若子类改变父类成员 offset，同一代码会访问错误位置。保持 offset 不变可保证父类代码对所有子类对象仍正确。

\subsection*{题 9：对 \texttt{a = a + 1}，为什么 \texttt{INC a} 可能比三条指令更好？}

\textbf{解答：} 三条指令实现需要 load、add、store，成本高且代码长。若目标机器支持 \texttt{INC a}，一条指令即可完成加一，通常更短、更快，指令成本也更低。

\subsection*{题 10：指令调度如何优化如下代码？}

\begin{verbatim}
A = x * y;
B = A + 1;
C = y;
\end{verbatim}

\textbf{解答：} 因为 \(B\) 依赖 \(A\)，而乘法可能有延迟，可把独立语句 \(C=y\) 移到中间：

\begin{verbatim}
A = x * y;
C = y;
B = A + 1;
\end{verbatim}

这样等待乘法结果时可以执行 \(C=y\)，减少空转。

\subsection*{题 11：寄存器干涉图中有边和无边分别表示什么？}

\textbf{解答：} 有边表示两个变量 live range 重叠，不能分配到同一寄存器；无边表示它们不会同时需要保留，可共享一个寄存器。k 可着色表示可用 k 个寄存器完成分配。

\subsection*{题 12：什么是寄存器溢出，为什么慢？}

\textbf{解答：} 当寄存器不够时，编译器把某些变量放到内存中，这叫 spilling。被溢出的变量使用前要从内存 load，用完后可能要 store 回内存。内存访问远慢于寄存器访问，因此溢出会显著降低性能。

\subsection*{题 13：什么是 jump to jump 窥孔优化？}

\textbf{解答：} 若条件跳转先跳到某标签，而该标签处又无条件跳到另一个标签，则可把条件跳转直接改到最终目标。例如：

\begin{verbatim}
if a < b goto L1
...
L1: goto L2
\end{verbatim}

优化为：

\begin{verbatim}
if a < b goto L2
...
L1: goto L2
\end{verbatim}

这样减少一次跳转。

\section*{复习建议}
\addcontentsline{toc}{section}{复习建议}

复习本讲建议抓住四条线：

\begin{itemize}
  \item 代码生成线：IR 到目标代码，重点是指令选择、寄存器分配、指令调度。
  \item 栈和寄存器线：栈式机器、累加器、MIPS 栈指针、保存恢复、表达式代码生成。
  \item 运行时线：内存布局、活动记录、\texttt{\$sp}/\texttt{\$fp}、调用约定、函数调用。
  \item 后端优化线：指令成本、调度、寄存器干涉图、图着色、溢出、窥孔优化。
\end{itemize}

最容易混淆的是 \texttt{\$sp} 和 \texttt{\$fp}：\texttt{\$sp} 管理栈顶，随压栈弹栈移动；\texttt{\$fp} 固定当前活动记录基准，适合稳定访问参数和局部变量。

\end{document}
